\chapter{Entorno socioeconómico}

En este capítulo, se presenta el entorno socioeconómico del proyecto, que comprende tanto el presupuesto para la elaboración de este Trabajo de Fin de Grado, como el impacto socio-económico del mismo en su sector.

\section{Presupuesto}

El objetivo de este apartado es detallar los aspectos económicos asociados al desarrollo del proyecto, dado que su realización conlleva una serie de costes. En este caso, se pueden dividir en dos categorías: costes de herramientas y costes humanos.

\subsubsection{Costes de herramientas}

Esta categoría incluye los recursos hardware y software que se han utilizado para el desarrollo de la aplicación y para la elaboración de la documentación. En cuanto al software, se ha hecho uso principalmente de herramientas gratuitas o con licencias proporcionadas por la universidad, por lo que este no ha supuesto un gasto real. No obstante, sí se considera el coste imputable del equipo utilizado, cuyo modelo y detalles se especificaron en el apartado de Medios empleados, situado al inicio de este documento.

Para calcular el coste imputable o amortizado del equipo, se utiliza esta fórmula:

\[
Coste amortizado = \left(\frac{m}{d}\right) \cdot C \cdot u
\]

Donde:
\begin{itemize}
	\item $m$: número de meses en los que el equipo se utiliza en el proyecto.
	\item $d$: número de meses del periodo de depreciación total o vida útil del equipo.
	\item $C$: coste del equipo sin IVA.
	\item $u$: porcentaje de uso dedicado al proyecto.
\end{itemize}

Teniendo como valores 5 meses de uso del equipo, 60 meses como periodo de depreciación, un coste de 951€, y un porcentaje de uso dedicado de un 100\%; nos queda un coste amortizado de 79,25€.

\begin{table}[H]
\centering
\begin{tabular}{|p{4.5cm}|p{7.5cm}|c|}
\hline
\textbf{Herramienta} & \textbf{Descripción} & \textbf{Coste (€)} \\
\hline
Ordenador portátil & Coste amortizado del equipo utilizado para desarrollo, pruebas y redacción de la memoria. & 79,25 \\
\hline
Lenguajes y frameworks de desarrollo & Tecnologías empleadas para frontend, backend y documentación (Latex) con licencia gratuita. & 0,00 \\
\hline
Visual Studio Code & Editor utilizado tanto para el código como para la documentación, con licencia gratuita. & 0,00 \\
\hline
Google Meet & Plataforma de comunicación para las reuniones de seguimiento con el tutor. & 0,00 \\
\hline
\textbf{TOTAL} &  & \textbf{79,25 €} \\
\hline
\end{tabular}
\caption{Costes de herramientas}
\end{table}

\subsubsection{Costes humanos}

El desarrollo del proyecto ha requerido una dedicación aproximada de 5 meses. Durante este periodo se distinguen dos roles principales: desarrollo y gestión del proyecto.

Por un lado, la autora del proyecto ha desempeñado el rol de desarrolladora (Ingeniera Junior), encargándose de la implementación completa del sistema. Por otro lado, también ha asumido el rol de Product Owner. Finalmente, el tutor ha actuado como Scrum Master.

A partir de la planificación realizada, se estima la siguiente dedicación:

\begin{table}[H]
\centering
\begin{tabular}{|c|c|c|c|c|c|}
\hline
\textbf{Perfil} & \textbf{Horas/mes} & \textbf{Meses} & \textbf{Horas totales} & \textbf{€/hora} & \textbf{Coste (€)} \\
\hline
Ingeniera Junior & 54 & 5 & 270 & 30 & 8.100,00 \\
\hline
Product Owner & 2 & 5 & 10 & 40 & 400,00 \\
\hline
Scrum Master & 4 & 5 & 20 & 40 & 800,00 \\
\hline
	extbf{TOTAL} &  &  & \textbf{316} &  & \textbf{9.300,00 €} \\
\hline
\end{tabular}
\caption{Costes humanos}
\end{table}

\subsubsection{Resumen de costes}

A partir de los costes anteriores, el presupuesto total del proyecto es el siguiente:

\begin{table}[H]
\centering
\begin{tabular}{|c|c|}
\hline
\textbf{Concepto} & \textbf{Coste (€)} \\
\hline
Costes de herramientas & 79,25 \\
\hline
Costes humanos & 9.300,00 \\
\hline
\textbf{Coste total (sin IVA)} & \textbf{9.379,25 €} \\
\hline
IVA (21\%) & 1.969,64 \\
\hline
\textbf{TOTAL (IVA incluido)} & \textbf{11.348,89 €} \\
\hline
\end{tabular}
\caption{Resumen del coste del proyecto}
\end{table}

En resumen, el desarrollo del proyecto presenta un coste concentrado en el esfuerzo humano que asciende a un total de 11.348,89 €.

\section{Impacto socioeconómico}

Hoy en día, el conocimiento de lenguas extranjeras se ha convertido en una competencia clave tanto a nivel académico como profesional. Es más, no son solo las instituciones académicas y empresas las entidades que optan por la internacionalización, sino que también la cultura se presenta como global en múltiples dominios como el cine, la música, el teatro, las redes sociales o el mundo de los videojuegos. Por ello, este proyecto busca poner el aprendizaje de idiomas al alcance de cualquiera con una conexión a Internet, ofreciendo una aplicación web con un enfoque personal, colaborativo y social. Así, el conjunto de sus beneficiarios incluye a cualquier persona interesada en mejorar sus habilidades lingüísticas por una amplia gama de motivos, desde académicos y profesionales hasta culturales, sociales y personales.

En primer lugar, la aplicación contribuye a crear un impacto en el \textbf{acceso a la educación en idiomas}, permitiendo a los usuarios aprender de forma libre y flexible ya que, a diferencia de los métodos tradicionales, una solución digital elimina barreras como la localización geográfica o los horarios, facilitando un aprendizaje continuo adaptado a los ritmos y disponibilidad de cada uno. La adopción de soluciones tecnológicas fomenta la innovación en el ámbito educativo, promoviendo métodos más interactivas, personalizadas y accesibles.

Además, la digitalización del aprendizaje contribuye a la \textbf{optimización de recursos}, lo que tiene un impacto positivo en el medio ambiente. Al tratarse de una aplicación software, se promueve un modelo más sostenible ya que se reducen los costes asociados a infraestructuras físicas, materiales educativos tradicionales, costes humanos adicionales y desplazamientos.

Otro impacto relevante es la mejora de la \textbf{empleabilidad} de los usuarios. En un mundo globalizado como el de hoy, el dominio de lenguas extranjeras es una habilidad cada vez más demandada en el mercado laboral, llegando a ser imprescindible en múltiples sectores como la consultoría, las ventas o la ingeniería. Esta aplicación puede contribuir al desarrollo y mejora de estas aptitudes, lo que aumenta las oportunidades laborales de los usuarios, especialmente en sectores internacionales o tecnológicos, que abarcan cada vez más espacio en el mercado.

Asimismo, la aplicación puede favorecer la \textbf{inclusión social} y la \textbf{integración cultural}. Por un lado, las competencias en idiomas permiten a las personas comunicarse en entornos diversos, facilitando la movilidad internacional, la integración de personas extranjeras y el intercambio cultural. Por otro lado, esta plataforma no solo contribuye a mejorar dichas competencias para sus usuarios, sino que es capaz de crear en si misma una comunidad que conecta gente de diferentes localidades, nacionalidades y contextos, y que, al estar enfocada a la redacción de temas libres en lugar de seguir un método rígido y cerrado, permite un intercambio real y personal de experiencias e ideas.

Desde el punto de vista económico, el desarrollo de una aplicación educativa no solo genera valor para los usuarios, sino que también pueden derivar en una \textbf{oportunidad de emprendimiento} con el desarrollo de un modelo de negocio digital. Esto, a su vez, contribuye al avance del cada vez más en auge sector "EdTech" o "Educational Technology", en español Tecnología Educativa, que está generando múltiples oportunidades de mercado.

En cuanto al impacto tecnológico, la realización de este proyecto ha implicado la adquisición y aplicación de conocimientos en áreas como el desarrollo de software, diseño de interfaces, programación de APIs y gestión de sistemas, lo que ha contribuido enormemente a mi formación en competencias altamente demandadas en el mercado laboral actual.

En resumen, el desarrollo de una aplicación para el aprendizaje de idiomas tiene un impacto socioeconómico positivo al mejorar la accesibilidad de la educación, aumentar la empleabilidad, favorecer la inclusión social, impulsar el sector tecnológico y contribuir a la transformación digital del aprendizaje. Estos beneficios repercuten tanto en los usuarios individuales como en la sociedad en su conjunto.